\section{Investment as a score intervention}
\label{sec:methods}
Let $X\in[0,1]^{m\times n}$ contain higher-is-better scores for $m$ evaluations on $n$ tasks. For target $t$, choose $S\subseteq\{1,\ldots,n\}$ and define
\begin{equation}
 x'_{tj}=x_{tj}+\mathbb{1}\{j\in S\}\alpha(u_j-x_{tj}),
 \label{eq:intervention}
\end{equation}
where $\alpha\in(0,1]$ and $u_j\geq\max_i x_{ij}$. All other rows remain fixed. We measure investment as $|S|$, also reporting $|S|/n$ within each benchmark. Equal task costs make a transparent combinatorial budget; subjects with different sizes and difficulties need not have equal training costs. Independent task improvements, no adverse cross-task effects, complete knowledge of the matrix and unconstrained subset selection are assumptions.

For aggregation rule $A$, let $A_i(X)$ be evaluation $i$'s aggregate. The minimum co-first budget is
\begin{equation}
 K_A(t;\alpha,u)=\min_S\{|S|: A_t(X')\geq\max_{i\ne t}A_i(X')\},
 \label{eq:budget}
\end{equation}
with $K_A=\infty$ if no subset succeeds. Existing winners have $K_A=0$. Co-first permits a numerical tolerance of $10^{-10}$; strict-first is analyzed separately. The main setting is $\alpha=1,u_j=1$. Sensitivities use $\alpha\in\{0.5,0.25,0.1\}$ at unit ceiling, and $\alpha=1,u_j=\max_i x_{ij}$. The observed-maximum ceiling limits gains to scores already represented in the snapshot; it still assumes the target can attain them.

\paragraph{Rules.}
We select four cardinal summaries to cover additive gains, central performance, multiplicative balance and weakest-task performance, plus one ordinal rule based on task ranks. This is a controlled recomputation of alternative rules on the same matrix, not a claim that each is deployed by the source leaderboard. We weight tasks equally. The mean is $n^{-1}\sum_j x_{ij}$; the geometric mean is $(\prod_j x_{ij})^{1/n}$, retaining zeros; the minimum is $\min_j x_{ij}$. The median is the central order statistic for odd $n$, or the average of the two central values for even $n$. For Borda, define $p(a,b)=\mathbb{1}\{a>b\}+\tfrac12\mathbb{1}\{a=b\}$ and
\begin{equation}
 B_i(X)=\sum_j\sum_{\ell\ne i}p(x_{ij},x_{\ell j}).
\end{equation}
This is the sum of average ascending task ranks minus one; higher is better. Its units differ from the cardinal scores, but $K_A$ has the same task-count units for every rule. The strongest rival can differ by rule, an essential part of the comparison.

\paragraph{Exact cardinal solutions.}
For a cardinal rule, rivals' aggregates remain fixed. Let $H$ denote the largest rival aggregate and $y_j=x_{tj}+\alpha(u_j-x_{tj})$ the allowed improved score. For the mean, select the shortest prefix of descending gains $y_j-x_{tj}$ whose sum reaches $nH-\sum_j x_{tj}$. Exchanging a selected smaller gain for an unselected larger one cannot hurt, proving optimality for unit costs. For a positive geometric threshold, first force upgrades of all zero coordinates; an unremovable zero implies infeasibility. Sort the remaining log gains and apply the same argument.

For $n=2h-1$, median success requires at least $h$ coordinates at least $H$. If $q$ already meet $H$, select $\max(0,h-q)$ eligible coordinates that can reach it; insufficient eligible coordinates imply infeasibility. Thus full unit-ceiling improvement requires at most $h$ selections. This argument does not apply to the conventional even median: all subsets are enumerated for the six-task benchmark. For the minimum, every coordinate below $H$ must reach $H$; their count is exact if all are eligible, and otherwise the target is infeasible. A low rival minimum can make this requirement inexpensive even when the target trails on many other metrics.

\paragraph{Borda with moving rival totals.}
Let $d_{ij}=p(y_j,x_{ij})-p(x_{tj},x_{ij})$ for $i\ne t$ and $g_j=\sum_{i\ne t}d_{ij}$. Selecting $j$ gains $g_j$ target points and removes $d_{ij}$ rival points. With binary $z_j$ the exact optimization is
\begin{align}
 \min_z\;&\sum_j z_j \quad \text{subject to}\nonumber\\
 &\sum_j(g_j+d_{ij})z_j\geq B_i(X)-B_t(X),\quad i\ne t.
 \label{eq:borda}
\end{align}
Freezing rivals' Borda totals would solve a different problem. We solve this binary program with zero relative optimality gap for HELM, and enumerate all subsets for Open LLM. Strict Borda first place raises the required gap by half a point, since scores are half-integral. Cardinal strict-first uses the stated numerical tolerance. Each feasible subset is checked directly against the winning condition.

For all five rules, improving target coordinates cannot worsen its relative position. Accordingly $K_A$ is nonincreasing in $\alpha$ and in the ceiling, treating infeasibility as infinity. This supplies a diagnostic, not an ordering between rules. Heterogeneous task costs would invalidate the sorted-gain count argument; coupled gains or changing rivals would change the feasible intervention set itself.
